\chapter{Kurzfassung}

\begin{german} %switch to German language rules

Freigängerkatzen bringen regelmäßig Beutetiere wie Mäuse oder Vögel in den Wohnraum. Herkömmliche Katzenklappen können dies nicht verhindern, da sie lediglich zwischen bekannten und fremden Tieren unterscheiden. Eine kamerabasierte Beuteerkennung ist schwierig, da das Beutetier im Bild klein und oft teilweise verdeckt ist und die Vorhersagen eines Modells durch Bewegungsunschärfe und wechselnde Lichtverhältnisse von Einzelbild zu Einzelbild schwanken. Entscheidungen auf Basis einzelner Bilder führen daher zu einem ständigen Wechsel zwischen Sperren und Freigeben der Klappe.

Die vorliegende Arbeit untersucht, ob eine Verarbeitungskette aus Objekterkennung, Objektverfolgung, Bildsegmentierung und zeitlicher Zustandsmodellierung eine robustere Beuteerkennung ermöglicht als die Auswertung einzelner Bilder. Ein feinabgestimmter YOLO-Objektdetektor lokalisiert die Gesichter der Katzen, ByteTrack ordnet die Erkennungen über die Zeit eindeutigen Identitäten zu, und ein YOLO-Segmentierungsmodell bewertet einen nach unten erweiterten Ausschnitt jedes Katzengesichts auf das Vorhandensein von Beute. Ein Zustandsautomat aggregiert die Konfidenzwerte jeder Katze über ein gleitendes Zeitfenster und entscheidet erst dann über die Freigabe der Klappe. Für den Nachtbetrieb arbeitet das System auf Graustufenbildern. Die Trainingsdaten des Segmentierungsmodells wurden mangels realer Aufnahmen mithilfe generativer KI erzeugt, und die Modelle wurden für den Betrieb auf einem Raspberry Pi 5 optimiert.

Die Evaluierung anhand von 20 KI-generierten Videosequenzen zeigt, dass das Gesamtsystem nur in 1,0\,\% der Einzelbilder mit Beute die Klappe fälschlicherweise freigibt und keine saubere Katze aussperrt. In Ablationsstudien senkt der Zustandsautomat den Sicherheitsfehler gegenüber der Einzelbildentscheidung von 17,0\,\% auf 1,0\,\% und den Komfortfehler von 19,7\,\% auf 3,3\,\%. Der zweistufige Ansatz sperrt saubere Katzen deutlich seltener als ein einstufiger Detektor (3,3\,\% gegenüber 29,7\,\%).

Die Ergebnisse bestätigen, dass die Kombination aus Objekterkennung, Objektverfolgung und zeitlicher Zustandsmodellierung deutlich stabilere Entscheidungen liefert als eine Einzelbildauswertung und sich auf kostengünstiger, lokaler Hardware umsetzen lässt.
\end{german}
